%% showhanmengrid.tex -- documentation for the showhanmengrid package
%%
%% Copyright (c) 2026 Munehiro Yamamoto <munepixyz@gmail.com>
%% License: LPPL 1.3c or later
%%
%% Build: make showhanmengrid.pdf
%%   (builds the example pages from the tests, then runs lualatex twice)
\documentclass{article}
\usepackage[margin=28mm]{geometry}
\usepackage{luatexja}
\usepackage{amsmath}
\usepackage{xcolor}
\usepackage{graphicx}
\usepackage[colorlinks, allcolors=blue!50!black]{hyperref}
\newcommand*\shgversion{v1.1 (2026/10/25)}
\newcommand*\cs[1]{\texttt{\char`\\#1}}
\newcommand*\pkg[1]{\textsf{#1}}
\newcommand*\key[1]{\texttt{#1}}
\newcommand*\meta[1]{\textlangle\textit{#1}\textrangle}
\newcommand\example[2]{%
  \begin{minipage}[t]{.48\linewidth}
    \centering
    \IfFileExists{#1}%
      {\fbox{\includegraphics[page=1,width=.9\linewidth]{#1}}}%
      {\fbox{\parbox[c][.3\textheight][c]{.9\linewidth}{\centering
         \ttfamily #1\\(run \texttt{make test} first)}}}%
    \par\smallskip{\small #2}
  \end{minipage}}

\title{The \pkg{showhanmengrid} package\\
  \large Show the hanmen grid of Japanese text and its overflow on every page}
\author{Munehiro Yamamoto\\
  \url{https://github.com/munepi/showhanmengrid}}
\date{\shgversion}

\begin{document}
\maketitle

\section{Introduction}

The \emph{hanmen} (kihon hanmen, 基本版面), the basic text block of
Japanese typesetting, is built from four values: the character size
(1\,zw, the width of a full-width character), the number of characters
per line, the number of lines per page and the line feed.  When
\begin{itemize}
\item \cs{textwidth} is a multiple of 1\,zw, and
\item \cs{textheight} is \cs{topskip} plus a multiple of
  \cs{baselineskip},
\end{itemize}
every character of the body text sits on a fixed grid: the text is
\emph{set solid} (\emph{beta-gumi}, ベタ組み).

The \pkg{showhanmengrid} package draws, in the background of every page,
\begin{enumerate}
\item \textcolor{cyan}{in cyan}, one square of 1\,zw $\times$ 1\,zw for
  each character position, and the frame of the whole grid;
\item \textcolor{magenta}{in magenta}, the \emph{overflow}: the part of
  the text area that does not fit the grid;
\item \textcolor{magenta}{in magenta}, when the character pitch is not
  1\,zw, a note in the top margin (section~\ref{sec:pitch}).
\end{enumerate}
Both horizontal (yoko, 横組み) and vertical (tate, 縦組み) writing are
supported.

\paragraph{Supported formats/engines}
pLaTeX and upLaTeX (with dvipdfmx), and LuaLaTeX with LuaTeX-ja.
LaTeX2e 2025-06-01 or later is required.  With any other engine
(pdfLaTeX, XeLaTeX, or LuaLaTeX without LuaTeX-ja) the package issues a
warning and does nothing.  Trim marks that move the page are taken into
account (section~\ref{sec:tombo}).

\section{Usage}

\begin{verbatim}
\usepackage{showhanmengrid}
\end{verbatim}
Nothing else is needed: the grid is computed at \verb|\begin{document}|
from the layout parameters of the class, and the writing direction is
taken from the document (\cs{tate} of pTeX, \texttt{direction=3} of
LuaTeX-ja).  Keys can be given as package options or with
\cs{showhanmengridsetup} in the preamble:
\begin{verbatim}
\usepackage[color=blue!50, overflow=false]{showhanmengrid}
\showhanmengridsetup{rulewidth=0.2pt}
\end{verbatim}

\begin{center}
\begin{tabular}{lll}
  \hline
  Key & Default & Description\\
  \hline
  \key{color}         & \texttt{cyan}    & color of the grid\\
  \key{overflowcolor} & \texttt{magenta} & color of the overflow\\
  \key{rulewidth}     & \texttt{0.1pt}   & width of the grid lines\\
  \key{frame}         & \texttt{true}    & draw the frame of the whole grid\\
  \key{overflow}      & \texttt{true}    & draw the overflow\\
  \key{chars}         & (computed)       & characters per line\\
  \key{lines}         & (computed)       & lines per page\\
  \key{ascent}        & \texttt{0.88zw}  & yoko: height of the character
                                           face above the baseline\\
  \key{pitch}         & \texttt{zw}      & pitch of the squares along the
                                           line: \texttt{zw},
                                           \texttt{natural} or a length\\
  \key{note}          & \texttt{true}    & draw the note on a character
                                           pitch other than 1\,zw\\
  \hline
\end{tabular}
\end{center}

\section{How the grid is computed}

Let 1\,zw and $b$ be the full-width character size and
\cs{baselineskip} of \cs{normalsize}, and $p$ the pitch of the squares
along the line (1\,zw by default; see section~\ref{sec:pitch}).  In
both yoko and tate, \cs{textwidth} is the length of a line and
\cs{textheight} the extent over which the lines are stacked, so that
\begin{align*}
  \text{chars} &= \lfloor (\cs{textwidth} - 1\,\text{zw}) / p \rfloor + 1,\\
  \text{lines} &= \lfloor (\cs{textheight} - \cs{topskip}) / b \rfloor + 1,
\end{align*}
and the grid is $(\text{chars}-1)\,p + 1\,\text{zw}$ long along the
lines (that is, $\text{chars}\times 1\,\text{zw}$ for $p=1\,\text{zw}$)
and $(\text{lines}-1)\,b + 1\,\text{zw}$ wide across the lines.
A remainder of less than 0.01\,pt is ignored (jlreq, for instance, adds
0.0001\,pt to \cs{textwidth}).  Both values and the overflow are written
to the log.

\paragraph{Yoko}
The grid begins at the top of the character face of the first line:
\cs{topskip} minus the ascent (0.88\,zw by default) below the top of the
text area.  The lines run from top to bottom.

\paragraph{Tate}
The baseline of the first line is \cs{topskip} left of the right edge of
the text area, and the character face extends 0.5\,zw on each side of
the baseline.  The lines run from right to left.

\paragraph{Side margins}
The text area of the current page is placed with \cs{oddsidemargin} or
\cs{evensidemargin} as the LaTeX kernel does it.  The \cs{@outputpage}
of pLaTeX has a branch that would swap the two margins in tate, but it
tests \cs{iftdir} inside \verb|\vbox{\yoko ...}| and so never applies;
the package follows what is actually typeset.

\section{The overflow}

\begin{description}
\item[Along the lines] $\cs{textwidth} - (\text{chars}-1)\,p -
  1\,\text{zw}$, drawn after the end of the lines (right in yoko, bottom
  in tate).  It is distributed into the inter-character spaces of every
  justified line: the characters drift away from their squares.
\item[Across the lines] $\cs{textheight} - \cs{topskip} -
  (\text{lines}-1)\,b$, drawn after the last line (bottom in yoko, left
  in tate).  The page is not filled by whole lines; with
  \cs{flushbottom} the remainder is distributed into vertical glue.
\end{description}

A text area with an overflow is what you get when only the four margins
are given, as in word processors or with the \pkg{geometry} package.
The default layout of the classic Japanese classes is another example:
jbook (pLaTeX), ujbook (upLaTeX) and ltjbook (LuaTeX-ja) take their
\cs{textwidth} from the standard book class,
$\min(\cs{paperwidth}-2\,\text{in}, 327\,\text{pt})$, a value chosen for
Computer Modern; at 10\,pt it is 33.98\,zw, i.e.\ 9.47\,pt longer than
33 characters.  Their \cs{textheight}, $\lfloor(\cs{paperheight}-3.5\,
\text{in})/b\rfloor\,b + \cs{topskip}$, is then truncated to whole
points by \cs{@settopoint}; since \cs{topskip} is not a whole number of
points, the truncation makes the last line no longer fit, and one line is
lost (39 lines instead of 40 on A4, with an overflow of about 14\,pt).
Only the \texttt{compatibility}-mode sizes of these classes use multiples
of \cs{Cwd}.  jsbook and ltjsbook set \cs{textwidth} to a multiple of
1\,zw, and jlreq sets the text solid by itself.

\begin{figure}[htbp]
\example{test-overflow-ujbook-uplatex.pdf}{ujbook, default layout (A4):
  33 chars $\times$ 39 lines with an overflow in both directions}\hfill
\example{test-jlreq-tate-uplatex.pdf}{jlreq, \texttt{tate}: 63 chars
  $\times$ 26 lines, set solid}
\end{figure}

\section{Character pitch other than \texorpdfstring{1\,zw}{1zw}}\label{sec:pitch}

Japanese text is not always set with the characters side by side.
\emph{Shibu-aki} (四分アキ, 0.25\,zw between characters) is common in
tate, and some designs spread 34 characters over the width of 35.  The
space can be given by \cs{kanjiskip}, which is simple, or, more
properly, by the glue between ideographs in the JFM (TFM).

The package measures the natural character pitch of \cs{normalsize} as
\[
  p_{\text{natural}} = \bigl(\text{width of 「漢漢漢」} -
  \text{width of 「漢」}\bigr) / 2,
\]
which includes both.  When it differs from 1\,zw, a one-line note is
drawn in the top margin of every page, just above the head, and written
to the log:
\begin{verbatim}
showhanmengrid: character pitch = 1.25zw, not 1zw
  (\kanjiskip = 0zw, JFM +0.25zw); squares at pitch=zw
\end{verbatim}
The parenthesis splits the extra space into the natural width of
\cs{kanjiskip} and the part given by the JFM.  (With LuaTeX-ja,
\cs{kanjiskip} of \cs{maxdimen} means that the value of the JFM is
used; the note then says ``\verb|\kanjiskip = JFM|''.)  The key
\key{note=false} turns the note off.

The key \key{pitch} sets the pitch $p$ of the squares along the line:
\begin{description}
\item[\texttt{pitch=zw}] (default) the squares of the hanmen design,
  1\,zw each and side by side.  The characters drift away from them.
\item[\texttt{pitch=natural}] the squares (still 1\,zw) are placed at
  $p_{\text{natural}}$, so that every character sits in its square; the
  number of characters per line and the overflow follow from $p$.
\item[\texttt{pitch=}\meta{length}] the pitch is given directly, e.g.\
  \texttt{pitch=1.25zw}.
\end{description}
The line feed is not affected.

\begin{figure}[htbp]
\example{test-jfm-natural-lualatex.pdf}{shibu-aki by the JFM,
  \texttt{pitch=natural}: 33 chars $\times$ 33 lines}\hfill
\example{test-jfm-lualatex.pdf}{the same text, \texttt{pitch=zw}: the
  grid of the design, 41 chars}
\end{figure}

\section{Trim marks and \texorpdfstring{\cs{mag}}{\textbackslash mag}}\label{sec:tombo}

The grid is drawn in the \texttt{shipout/background} picture, whose
origin moves with \cs{hoffset} and \cs{voffset} as the text does.  In
addition,
\begin{itemize}
\item with pLaTeX, the \texttt{tombo} option of jclasses and jsclasses,
  and the classes and packages that rely on it (\cs{iftombow} is true;
  e.g.\ gentombow and jlreq-trimmarks on pLaTeX), move the page body by
  1\,in in both directions in \cs{@outputpage}; the grid is moved
  likewise;
\item with LuaLaTeX, gentombow and jlreq-trimmarks move the page with
  \cs{hoffset} and \cs{voffset}, which needs no correction.
\end{itemize}
A \cs{mag} other than 1000 is not supported: the package writes a
warning to the log and draws nothing.

\section{Limitations}

\begin{itemize}
\item Two-column layout is not supported: the grid spans \cs{textwidth}.
\item \cs{mag} other than 1000 is not supported (see above).
\item The character pitch is measured once, with 「漢」 of
  \cs{normalsize}; a JFM that gives different classes of characters
  different spaces is represented by the ideographs.
\item The grid is computed once, at \verb|\begin{document}|; changes of
  \cs{textwidth} or \cs{textheight} later in the document are ignored.
\end{itemize}

\section{Tests}

\verb|make test| compiles every test file with pLaTeX, upLaTeX
(+ dvipdfmx) and LuaLaTeX and checks the log.  \texttt{SHG-EXPECT OK}
means that the writing direction, the number of characters and lines,
and the overflow are as expected; \texttt{SHG-POS OK} means that the
first character of the page, measured with \cs{pdfsavepos} (pTeX) or
\cs{savepos} (LuaTeX), sits in the top left square of the grid as drawn,
on odd and even pages; some tests also check the end of the last square
of a line.  \texttt{SHG-NOTE OK} means that the note on the character
pitch is as expected (or absent).  With LuaTeX-ja in tate the position
check is skipped
(\texttt{SHG-POS SKIP}): LuaTeX-ja typesets tate text in a rotated box
whose \cs{savepos} coordinates are not those of the page.

\section{History}

\begin{description}
\item[v1.1 (2026/10/25)] Character pitch other than 1\,zw: the note in
  the top margin and the keys \key{pitch} and \key{note}.  Trim marks
  (\cs{iftombow} of pLaTeX, gentombow, jlreq-trimmarks) are taken into
  account.  \cs{mag} other than 1000 disables the package with a
  warning.  Fixed: with a non-zero \cs{hoffset} or \cs{voffset}, the
  grid was shifted by that amount twice.
\item[v1.0 (2026/10/25)] First public release.  Generalized from the
  version written for the book on LaTeX page layout: it no longer needs
  \cs{hanmenwidth}/\cs{hanmenheight}, supports tate and pLaTeX, and
  draws the overflow.
\end{description}

\section{License}

This package is licensed under the conditions of the \LaTeX{} Project
Public License, either version 1.3c of this license or any later version.

Copyright \copyright{} 2026 Munehiro Yamamoto

\end{document}
